% Aufgaben für den LaTeX-Kurs (2015-08-07)
% Autor: Simon May
% Datum: 2015-08-07

\documentclass[a4paper, 12pt, headsepline, pagesize, parskip=half-, headlines=3.1]{scrartcl}

\usepackage{fontspec}
\defaultfontfeatures{Ligatures=TeX}
\usepackage{polyglossia}
\setdefaultlanguage{german}

\usepackage[intlimits]{amsmath}
\usepackage{newunicodechar}
% Autor: Simon May
% Datum: 2015-03-04

% Diese Datei ermöglicht es, Mathe-Symbole (z.B. \gamma) direkt als
% Sonderzeichen (d.h. γ) einzugeben

% silence unterdrückt Warnungen; vor hyperref laden
\usepackage{silence}
\WarningFilter[pdflatex-unicode-math]{newunicodechar}{Redefining Unicode character}
\ActivateWarningFilters[pdflatex-unicode-math]
\WarningFilter{microtype}{Unknown slot number of character}
\AtBeginDocument{\DeactivateWarningFilters}

\newunicodechar{†}{\dag}
\newunicodechar{‡}{\ddag}
\newunicodechar{…}{\ldots}
\newunicodechar{⋯}{\cdots}
\newunicodechar{⋮}{\vdots}
\newunicodechar{⋱}{\ddots}
\newunicodechar{⋰}{\iddots}
\newunicodechar{α}{\alpha}
\newunicodechar{β}{\beta}
\newunicodechar{γ}{\gamma}
\newunicodechar{δ}{\delta}
\newunicodechar{ε}{\varepsilon}
\newunicodechar{ϵ}{\epsilon}
\newunicodechar{ζ}{\zeta}
\newunicodechar{η}{\eta}
\newunicodechar{θ}{\theta}
\newunicodechar{ϑ}{\vartheta}
\newunicodechar{ι}{\iota}
\newunicodechar{κ}{\kappa}
\newunicodechar{ϰ}{\varkappa}
\newunicodechar{λ}{\lambda}
\newunicodechar{μ}{\mu}
\newunicodechar{ν}{\nu}
\newunicodechar{ξ}{\xi}
\newunicodechar{ο}{o}
\newunicodechar{π}{\pi}
\newunicodechar{ρ}{\rho}
\newunicodechar{ϱ}{\varrho}
\newunicodechar{σ}{\sigma}
\newunicodechar{τ}{\tau}
\newunicodechar{υ}{\upsilon}
\newunicodechar{φ}{\varphi}
\newunicodechar{ϕ}{\phi}
\newunicodechar{χ}{\chi}
\newunicodechar{ψ}{\psi}
\newunicodechar{ω}{\omega}
\newunicodechar{Α}{\mathrm{A}}
\newunicodechar{Β}{\mathrm{B}}
\newunicodechar{Γ}{\Gamma}
\newunicodechar{Δ}{\Delta}
\newunicodechar{Ε}{\mathrm{E}}
\newunicodechar{Ζ}{\mathrm{Z}}
\newunicodechar{Η}{\mathrm{H}}
\newunicodechar{Θ}{\Theta}
\newunicodechar{Ι}{\mathrm{I}}
\newunicodechar{Κ}{\mathrm{K}}
\newunicodechar{Λ}{\Lambda}
\newunicodechar{Μ}{\mathrm{M}}
\newunicodechar{Ν}{\mathrm{N}}
\newunicodechar{Ξ}{\Xi}
\newunicodechar{Ο}{\mathrm{O}}
\newunicodechar{Π}{\Pi}
\newunicodechar{Ρ}{\mathrm{P}}
\newunicodechar{Σ}{\Sigma}
\newunicodechar{Τ}{\mathrm{T}}
\newunicodechar{Υ}{\Upsilon}
\newunicodechar{Φ}{\Phi}
\newunicodechar{Χ}{\Chi}
\newunicodechar{Ψ}{\Psi}
\newunicodechar{Ω}{\Omega}
\newunicodechar{∑}{\sum}
\newunicodechar{∫}{\int}
\newunicodechar{∬}{\iint}
\newunicodechar{∭}{\iiint}
\newunicodechar{⨌}{\iiiint}
\newunicodechar{∮}{\oint}
\newunicodechar{∯}{\oiint}
\newunicodechar{∰}{\oiiint}
\newunicodechar{∇}{\nabla}
\newunicodechar{∂}{\partial}
\newunicodechar{√}{\sqrt}
\newunicodechar{∈}{\in}
\newunicodechar{∋}{\ni}
\newunicodechar{∉}{\notin}
\newunicodechar{∀}{\forall}
\newunicodechar{∃}{\exists}
\newunicodechar{∄}{\nexists}
\newunicodechar{∴}{\therefore}
\newunicodechar{∵}{\because}
\newunicodechar{〈}{\langle}
\newunicodechar{〉}{\rangle}
\newunicodechar{⌊}{\lfloor}
\newunicodechar{⌋}{\rfloor}
\newunicodechar{⌈}{\lceil}
\newunicodechar{⌉}{\rceil}
\newunicodechar{∼}{\sim}
\newunicodechar{∝}{\propto}
\newunicodechar{∞}{\infty}
\newunicodechar{ℵ}{\aleph}
\newunicodechar{ℏ}{\hbar}
\newunicodechar{℘}{\wp}
\newunicodechar{ℓ}{\ell}
\newunicodechar{∅}{\emptyset}
\newunicodechar{×}{\times}
\newunicodechar{⋅}{\cdot}
\newunicodechar{÷}{\div}
\newunicodechar{⋆}{\star}
\newunicodechar{∘}{\circ}
\newunicodechar{⋄}{\diamond}
\newunicodechar{⊕}{\oplus}
\newunicodechar{⊖}{\ominus}
\newunicodechar{⊗}{\otimes}
\newunicodechar{⊘}{\oslash}
\newunicodechar{⊙}{\odot}
\newunicodechar{±}{\pm}
\newunicodechar{∓}{\mp}
\newunicodechar{≈}{\approx}
\newunicodechar{≡}{\equiv}
\newunicodechar{≠}{\ne}
\newunicodechar{≥}{\ge}
\newunicodechar{≤}{\le}
\newunicodechar{≫}{\gg}
\newunicodechar{≪}{\ll}
\newunicodechar{⊂}{\subset}
\newunicodechar{⊃}{\supset}
\newunicodechar{⊆}{\subseteq}
\newunicodechar{⊇}{\supseteq}
\newunicodechar{⊈}{\nsubseteq}
\newunicodechar{⊉}{\nsupseteq}
\newunicodechar{≔}{\coloneqq}
\newunicodechar{≕}{\eqqcolon}
\newunicodechar{¬}{\neg}
\newunicodechar{∨}{\vee}
\newunicodechar{∧}{\wedge}
\newunicodechar{∪}{\cup}
\newunicodechar{∩}{\cap}
\newunicodechar{⋁}{\bigvee}
\newunicodechar{⋀}{\bigwedge}
\newunicodechar{⋃}{\bigcup}
\newunicodechar{⋂}{\bigcap}
\newunicodechar{⟂}{\perp}
\newunicodechar{∥}{\parallel}
\newunicodechar{∦}{\nparallel}
\newunicodechar{𝚤}{\imath}
\newunicodechar{𝚥}{\jmath}
\newunicodechar{⇔}{\Leftrightarrow}
\newunicodechar{⇕}{\Updownarrow}
\newunicodechar{⇐}{\Leftarrow}
\newunicodechar{⇒}{\Rightarrow}
\newunicodechar{⇑}{\Uparrow}
\newunicodechar{⇓}{\Downarrow}
\newunicodechar{↔}{\leftrightarrow}
\newunicodechar{↕}{\updownarrow}
\newunicodechar{←}{\leftarrow}
\newunicodechar{→}{\rightarrow}
\newunicodechar{↑}{\uparrow}
\newunicodechar{↓}{\downarrow}
\newunicodechar{⟷}{\longleftrightarrow}
\newunicodechar{⟵}{\longleftarrow}
\newunicodechar{⟶}{\longrightarrow}
\newunicodechar{⇇}{\leftleftarrows}
\newunicodechar{⇉}{\rightrightarrows}
\newunicodechar{⇈}{\upuparrows}
\newunicodechar{⇊}{\downdownarrows}
\newunicodechar{⟺}{\Longleftrightarrow}
\newunicodechar{⟸}{\Longleftarrow}
\newunicodechar{⟹}{\Longrightarrow}
\newunicodechar{↦}{\mapsto}
\newunicodechar{↤}{\mapsfrom}
\newunicodechar{⟼}{\longmapsto}
\newunicodechar{⟻}{\longmapsfrom}
\newunicodechar{⟾}{\Longmapsto}
\newunicodechar{⟽}{\Longmapsfrom}
\newunicodechar{↗}{\nearrow}
\newunicodechar{↖}{\nwarrow}
\newunicodechar{↘}{\searrow}
\newunicodechar{↙}{\swarrow}
\newunicodechar{↩}{\hookleftarrow}
\newunicodechar{↪}{\hookrightarrow}
\newunicodechar{↶}{\curvearrowleft}
\newunicodechar{↷}{\curvearrowright}
\newunicodechar{↺}{\circlearrowleft}
\newunicodechar{↻}{\circlearrowright}
\newunicodechar{↫}{\looparrowleft}
\newunicodechar{↬}{\looparrowright}
\newunicodechar{⇋}{\leftrightharpoons}
\newunicodechar{⇌}{\rightleftharpoons}
\newunicodechar{↼}{\leftharpoonup}
\newunicodechar{↽}{\leftharpoondown}
\newunicodechar{⇀}{\rightharpoonup}
\newunicodechar{⇁}{\rightharpoondown}
\newunicodechar{↿}{\upharpoonleft}
\newunicodechar{↾}{\upharpoonright}
\newunicodechar{⇃}{\downharpoonleft}
\newunicodechar{⇂}{\downharpoonright}
\newunicodechar{𝔸}{\mathbb{A}}
\newunicodechar{𝔹}{\mathbb{B}}
\newunicodechar{ℂ}{\mathbb{C}}
\newunicodechar{𝔻}{\mathbb{D}}
\newunicodechar{𝔼}{\mathbb{E}}
\newunicodechar{𝔽}{\mathbb{F}}
\newunicodechar{𝔾}{\mathbb{G}}
\newunicodechar{ℍ}{\mathbb{H}}
\newunicodechar{𝕀}{\mathbb{I}}
\newunicodechar{𝕁}{\mathbb{J}}
\newunicodechar{𝕂}{\mathbb{K}}
\newunicodechar{𝕃}{\mathbb{L}}
\newunicodechar{𝕄}{\mathbb{M}}
\newunicodechar{ℕ}{\mathbb{N}}
\newunicodechar{𝕆}{\mathbb{O}}
\newunicodechar{ℙ}{\mathbb{P}}
\newunicodechar{ℚ}{\mathbb{Q}}
\newunicodechar{ℝ}{\mathbb{R}}
\newunicodechar{𝕊}{\mathbb{S}}
\newunicodechar{𝕋}{\mathbb{T}}
\newunicodechar{𝕌}{\mathbb{U}}
\newunicodechar{𝕍}{\mathbb{V}}
\newunicodechar{𝕎}{\mathbb{W}}
\newunicodechar{𝕏}{\mathbb{X}}
\newunicodechar{𝕐}{\mathbb{Y}}
\newunicodechar{ℤ}{\mathbb{Z}}
\newunicodechar{𝒜}{\mathcal{A}}
\newunicodechar{ℬ}{\mathcal{B}}
\newunicodechar{𝒞}{\mathcal{C}}
\newunicodechar{𝒟}{\mathcal{D}}
\newunicodechar{ℰ}{\mathcal{E}}
\newunicodechar{ℱ}{\mathcal{F}}
\newunicodechar{𝒢}{\mathcal{G}}
\newunicodechar{ℋ}{\mathcal{H}}
\newunicodechar{ℐ}{\mathcal{I}}
\newunicodechar{𝒥}{\mathcal{J}}
\newunicodechar{𝒦}{\mathcal{K}}
\newunicodechar{ℒ}{\mathcal{L}}
\newunicodechar{ℳ}{\mathcal{M}}
\newunicodechar{𝒩}{\mathcal{N}}
\newunicodechar{𝒪}{\mathcal{O}}
\newunicodechar{𝒫}{\mathcal{P}}
\newunicodechar{𝒬}{\mathcal{Q}}
\newunicodechar{ℛ}{\mathcal{R}}
\newunicodechar{𝒮}{\mathcal{S}}
\newunicodechar{𝒯}{\mathcal{T}}
\newunicodechar{𝒰}{\mathcal{U}}
\newunicodechar{𝒱}{\mathcal{V}}
\newunicodechar{𝒲}{\mathcal{W}}
\newunicodechar{𝒳}{\mathcal{X}}
\newunicodechar{𝒴}{\mathcal{Y}}
\newunicodechar{𝒵}{\mathcal{Z}}
\newunicodechar{𝕬}{\mathfrak{A}}
\newunicodechar{𝕭}{\mathfrak{B}}
\newunicodechar{𝕮}{\mathfrak{C}}
\newunicodechar{𝕯}{\mathfrak{D}}
\newunicodechar{𝕰}{\mathfrak{E}}
\newunicodechar{𝕱}{\mathfrak{F}}
\newunicodechar{𝕲}{\mathfrak{G}}
\newunicodechar{𝕳}{\mathfrak{H}}
\newunicodechar{𝕴}{\mathfrak{I}}
\newunicodechar{𝕵}{\mathfrak{J}}
\newunicodechar{𝕶}{\mathfrak{K}}
\newunicodechar{𝕷}{\mathfrak{L}}
\newunicodechar{𝕸}{\mathfrak{M}}
\newunicodechar{𝕹}{\mathfrak{N}}
\newunicodechar{𝕺}{\mathfrak{O}}
\newunicodechar{𝕻}{\mathfrak{P}}
\newunicodechar{𝕼}{\mathfrak{Q}}
\newunicodechar{𝕽}{\mathfrak{R}}
\newunicodechar{𝕾}{\mathfrak{S}}
\newunicodechar{𝕿}{\mathfrak{T}}
\newunicodechar{𝖀}{\mathfrak{U}}
\newunicodechar{𝖁}{\mathfrak{V}}
\newunicodechar{𝖂}{\mathfrak{W}}
\newunicodechar{𝖃}{\mathfrak{X}}
\newunicodechar{𝖄}{\mathfrak{Y}}
\newunicodechar{𝖅}{\mathfrak{Z}}

\DeactivateWarningFilters[pdflatex-unicode-math]


\usepackage{selnolig}
\usepackage{microtype}
\usepackage{xcolor}
\usepackage{scrlayer-scrpage}
\usepackage{lastpage}
\usepackage{datetime}
\usepackage{csquotes}
\usepackage{siunitx}
\usepackage{enumitem}
\usepackage{wrapfig}
\usepackage{float}
\usepackage{listingsutf8}
\usepackage[backend=biber, style=verbose]{biblatex}
\usepackage{xfrac}
\usepackage{subcaption}

\usepackage[unicode]{hyperref}
\usepackage{cleveref}

% --- Paket-Einstellungen
% -- latex
\widowpenalty=9000
\clubpenalty=9000
% -- koma
\setkomafont{pageheadfoot}{\normalfont}
\addtokomafont{pagefoot}{\small}
% -- scrlayer-scrpage
\ihead[]{\textbf{Einführung in \LaTeX\\Aufgaben}}
\chead[]{\includegraphics[height=2.8\baselineskip]{res/logo.pdf}}
\ohead[]{Simon May\\Fachschaft Physik}
\cfoot[Seite: {\pagemark} von \pageref{LastPage}]{Seite: \pagemark\ von \pageref{LastPage}}
\pagestyle{scrheadings}
% -- biblatex
\addbibresource{res/literatur.bib}
% -- csquotes
\MakeOuterQuote{"}
% -- siunitx
\sisetup{
	locale=DE,
	output-product=\cdot,
	separate-uncertainty,
	quotient-mode=fraction,
	per-mode=fraction,
	fraction-function=\sfrac,
	math-micro=\text{µ},
	text-micro=µ
}
% -- listings
\colorlet{keywordcolor}{purple!60!black}
\colorlet{commentcolor}{green!40!black}
\colorlet{texcsicolor}{blue}
\colorlet{texcsiicolor}{blue!50!black}
\colorlet{bracescolor}{orange!70!black}
\colorlet{dollarcolor}{red}
\lstdefinestyle{mystandard}{
	% Stelle UTF8-Input mit Latin1 dar
	%inputencoding=utf8/latin1,
	% Zeilennummern links
	%numbers=left,
	% Darstellung der Zeilennummern
	numberstyle=\ttfamily,
	% Rahmen (links)
	%frame=L,
	columns=flexible,
	%xleftmargin=0.5cm,
	% Default Schriftart/Größe
	basicstyle=\ttfamily,
	% Darstellung für Keywords
	keywordstyle=\color{keywordcolor},
	% Darstellung für Kommentare
	commentstyle=\color{commentcolor},
	% Darstellung für Strings
	stringstyle=\color[RGB]{ 47,   0, 255},
	% Darstellung für TeX-Befehle; Stern: mit backslash
	texcsstyle=*\color{texcsicolor},
	texcsstyle=*[2]\color{texcsiicolor},
	% Tab-Breite
	tabsize=4,
	% Leerzeichen in Strings hervorheben
	showstringspaces=false,
	% Tabs hervorheben
	showtabs=false,
	% Leerzeichen hervorheben
	showspaces=false,
	% Alle Leerzeichen behalten
	keepspaces=true,
	% "Normale" Anführungszeichen: ""
	upquote=true,
	% Lange Zeilen aufbrechen
	breaklines=true,
	% Zeilen an Whitespace aufbrechen
	breakatwhitespace=true,
}

\lstset{
	language=[LaTeX]{TeX},
	style=mystandard,
	escapechar=|,
	deletetexcs={begin, end, documentclass, usepackage},
	moretexcs=[2]{begin, end, documentclass, usepackage},
	moretexcs={
		f, latex, Alpha,
		part, chapter, subsection, subsubsection, paragraph, subparagraph, color, text, mathbb, mathcal, tableofcontents, includegraphics, cref, crefrange, autocite, footcite, printbibliography,
		SI, num, si, SIrange, SIlist, m, per, s, squared, cm, nm, MHz, mega, Hz, cubed, tothe, mA, kV
	},
	literate=
		{\$}{\textcolor{dollarcolor}{\$}}1
		{\{}{\textcolor{bracescolor}{\{}}1
		{\}}{\textcolor{bracescolor}{\}}}1,
}


% --- Befehle
\newcommand{\email}[1]{\href{mailto:#1}{\texttt{#1}}}

\newcounter{aufgabe}
\newenvironment{aufgabe}[1][]{\stepcounter{aufgabe}
\noindent\textbf{Aufgabe \arabic{aufgabe}.} #1\par
\begin{enumerate}[label=\alph*)]}{\end{enumerate}\bigskip}
\newenvironment{aufgabe*}[1][]{\stepcounter{aufgabe}
\noindent\textbf{Aufgabe \arabic{aufgabe}.} #1}{\bigskip}

\begin{document}

\title{Einführung in \LaTeX: Aufgaben}
\author{Simon May\\\email{simon.may@uni-muenster.de}}
\date{\formatdate{7}{8}{2015}}

\maketitle

\begin{aufgabe}[Textformatierung]
	\item Stellt das folgende Beispiel aus der Vorlesung nach:
	\begin{itemize}
		\item \textbf{fett}
		\item \textit{kursiv}
		\item \texttt{Schreibmaschine}
		\item \underline{unterstrichen}
		\item \textsc{Kapitälchen}
	\end{itemize}
	\item Stellt diesen Satz nach:
	
	{\huge"\tiny Das ist \small ein Satz, \large der \Large diese \huge Schriftgrößen demonstriert."}
\end{aufgabe}

\begin{aufgabe*}
	Aufzählungen – Stellt das folgende Beispiel nach:
	\begin{enumerate}
		\item Physik~I
		\begin{itemize}
			\item Methodik der Physik
			\item Klassische (newtonsche) Mechanik
			\item Bewegungsgleichungen
			\item Einblick in die spezielle Relativitätstheorie
		\end{itemize}
		\item Physik~II
		\begin{enumerate}
			\item Kinetische Gastheorie
			\begin{enumerate}
				\item Maxwell-Boltzmann-Verteilung
				\item Ideales Gas
				\item Modelle realer Gase
			\end{enumerate}
			\item Klassische Thermodynamik
			\item Elektrostatik
			\item Magnetostatik
		\end{enumerate}
		\item Physik~III
		\begin{description}
			\item[Klassische Elektrodynamik] Maxwell-Gleichungen, Wechselstrom, Licht
			\begin{enumerate}
				\item Elektrodynamik im Vakuum
				\item Elektrodynamik in Materie
			\end{enumerate}
			\item[Ausbreitung von Wellen] Wellen auf Kabeln, Lichtausbreitung im Vakuum
			\item[Quantenphysik] Einblick in die Quantenmechanik
			\item[Spezielle Relativitätstheorie] Schnell bewegte Bezugssysteme
		\end{description}
	\end{enumerate}
\end{aufgabe*}

\begin{aufgabe}[Formeln]
	\item Setzt die folgenden Formeln mit \LaTeX:
	\begin{equation}
		E = mc^2
	\end{equation}	
	\begin{equation}
		\label{eq:newton}
		F = -G ⋅ \frac{mM}{r^2}
	\end{equation}	
	\begin{equation}
		C = \frac{Q}{U}
	\end{equation}	
	\begin{equation}
		U = U_0 ⋅ e^{-\frac{t}{RC}}
	\end{equation}	
	\begin{equation}
		\label{eq:circle_circumference}
		U = 2πr
	\end{equation}
	\begin{equation}\label{eq:circle_area}
		A = πr^2 = π\frac{d^2}{4}
	\end{equation}
	\begin{equation}
		r = \sqrt{\frac{(s-a)(s-b)(s-c)}{s}}
	\end{equation}
	\item Fügt eine Formel in den Fließtext ein:
	
	"Der Satz von Pythagoras, $a^2 + b^2 = c^2$, beschreibt das Verhältnis von Längen in einem rechtwinkligen Dreieck und ist eine sehr bekannte mathematische Aussage."
	\item Setzt folgende Herleitung mit \LaTeX:
	\begin{align*}
		(a + b)^3
		&= (a + b)^2 (a + b) \\
		&= (a^2 + 2ab + b^2) (a + b) \\
		&= \left[(a^3 + 2 a^2 b + a b^2) + (a^2 b + 2 a b^2 + b^3)\right] \\
		&= a^3 + 3 a^2 b + 3 a b^2 + b^3
	\end{align*}
	\item Versucht euch als Nächstes an diesen Formeln:
	\begin{align*}
		0 &= \frac{1}{π} ∫_0^{2π} \cos(nx) \sin(mx) \, \mathrm{d}x\\
		\dot{Q} &= - λA \frac{ΔT}{L}\\
		I &= ∫_A \vec{j} ⋅ \mathrm{d}\vec{A}\\
		\vec{F} &= - \frac{q_1 q_2}{4πε_0} \frac{\vec{r}_1 - \vec{r}_2}{|\vec{r}_1 - \vec{r}_2|^3} \\
		A &= \frac{\tan(x) \sin(x)}{π^2 ρ_{\text{neu}}} ⋅ ρ_{\text{alt}} ⋅ ∫_1^X \exp\!\left(\frac{m v^2}{2\ln(x)}\right) \, \mathrm{d}x
	\end{align*}
	\item Setzt die Maxwell-Gleichungen in Materie mit \LaTeX:
	\begin{align}
		\label{eq:max1}
		∇ ⋅ \vec{D} &= ρ_\text{f} \\
		∇ ⋅ \vec{B} &= 0 \\
		∇ × \vec{E} &= -\frac{∂\vec{B}}{∂t} \\
		\label{eq:max4}
		∇ × \vec{H} &= \vec{j}_\text{f} + \frac{∂\vec{D}}{∂t}
	\end{align}
	
	\emph{Erinnerung:} \lstinline!\nabla! ergibt $∇$, \lstinline!\partial! ergibt $∂$.
\end{aufgabe}

\clearpage

\begin{aufgabe}[Verweise]
	\item Versucht, einen automatischen Verweis auf \cref{eq:newton} zu erzeugen.
	\item Erzeugt den folgenden Text mit automatischen Verweisen:\\
	"\cref{eq:newton,eq:circle_circumference,eq:circle_area}".
	\item Erzeugt den folgenden Text mit automatischen Verweisen:\\
	"\crefrange{eq:max1}{eq:max4}".
	
	\emph{Hinweis:} Einen Bereich von Verweisen kann man mit\\
	\lstinline!\crefrange{<label1>}{<label2>}! erzeugen.
\end{aufgabe}

\begin{aufgabe}[Tabellen – Reproduziert die folgenden Tabellen.]
	\item 
	\begin{tabular}{|c c || c|}
		\hline
		1 & 2 & 3 \\ \hline\hline
		4 & 5 & 6 \\
		7 & 8 & 9 \\ \hline
	\end{tabular}
	\item
	\begin{tabular}{l | S S S S | S}
		& \multicolumn{4}{c|}{\textbf{Augenfarbe}} & \\
		\textbf{Haarfarbe}	& \textbf{Braun}	& \textbf{Blau} & \textbf{Grau}
			& \textbf{Grün}		& \textbf{Summe} \\ \hline
		\textbf{Braun}		& 119	& 84	& 54	& 29	& 286	\\
		\textbf{Blond}		& 7		& 94	& 10	& 16	& 127	\\
		\textbf{Schwarz}	& 68	& 20	& 15	& 5		& 108	\\
		\textbf{Rot}		& 26	& 17	& 14	& 14	& 71	\\ \hline
		\textbf{Summe}		& 220	& 215	& 93	& 64	& 592
	\end{tabular}
	
	\emph{Hinweis:} Mit dem Befehl\\
	\lstinline!\multicolumn{<spaltenzahl>}{<neue_spalte>}{Text...}! lassen sich Spalten zusammenfassen.

	\emph{Hinweis:} Zur einfacheren Eingabe gibt es z.\,B.\ in TeXStudio die Option "Wizards $→$ Quick Tabular…"
	\item Cyclohexan:
	
	\begin{tabular}{S S}
		\hline
		{Masse / \si{\g}}	& {Volumenänderung / \si{\ml}} \\
		\hline
		0,15		& 55+-2 \\
		0,18		& 59+-1 \\
		0,19		& 58+-1 \\
		0,19		& 55+-1 \\
		0,2			& 61+-3 \\
		\hline
	\end{tabular}
		
	\emph{Hinweis:} Verwendet \lstinline!S!-Spalten!
	
	\clearpage
	\item ~
	
	\begin{table}[H]
		\centering
		\caption{Gemessene Massen und daraus resultierende Volumenänderung bei der
			Bestimmung der molaren Masse von Ethanol mit der Dampfdichtemethode.}
		\begin{tabular}{S S}
			\hline
			{Masse / \si{\g}}	& {Volumenänderung / \si{\ml}} \\
			\hline
			0,06		& 55+-3 \\
			0,1			& 60+-2 \\
			0,11		& 58+-1 \\
			0,07		& 54+-1 \\
			0,08		& 52+-1 \\
			\hline
		\end{tabular}
		\label{tab:ethanol}
	\end{table}
\end{aufgabe}

\begin{aufgabe}[Abbildungen]
	\newcommand{\kursangle}{240}
	\item Fügt folgendes Bild (\texttt{answer.jpg}) in euer Dokument ein:
	
	\includegraphics[width=5cm, angle=\kursangle]{res/answer.jpg}
	
	\emph{Hinweis:} Der Winkel ist \SI{\kursangle}{\degree}.
	\item Erstellt die folgende Abbildung mit Unterabbildungen\\
	(zweites Bild: \texttt{8bit\_mario.png}):
	
	\begin{figure}[H]
		\centering
		\begin{subfigure}{0.45\textwidth}
		\centering
		\includegraphics[width=0.65\textwidth]{res/answer.jpg}
		\caption{Antwort auf die ultimative Frage nach dem Leben, dem Universum und dem ganzen Rest.}
		\end{subfigure}
		\hfill
		\begin{subfigure}{0.45\textwidth}
		\centering
		\includegraphics[width=0.65\textwidth]{res/8bit_mario.png}
		\caption{8bit-Mario sucht die Prinzessin.}
		\end{subfigure}
		\caption{Zwei Beispielbilder.}
	\end{figure}
	\item Erzeugt den folgenden Text mit eingebettetem Bild (\texttt{logo.pdf}):
\end{aufgabe}


\enlargethispage{\baselineskip}

Passend zum Internationalen Jahr des Lichts hat die Fachschaft Physik mit Herrn Prof.\ Dr.\ Michael Rohlfing den Lehrenden ausgezeichnet, der im vergangenen Jahr den Studierenden Elektromagnetismus, Optik und die spezielle Relativitätstheorie, also die Physik des Lichts, näher bringen konnte.

\begin{wrapfigure}{r}{0cm}
	\centering
	\includegraphics[width=9cm]{res/logo.pdf}
	\caption{Das Logo der Fachschaft Physik der Westfälischen Wilhelms-Universität Münster.}
\end{wrapfigure}
Das Lehrpreiskomitee hatte an der Spitze mit viel Licht zu kämpfen und musste zwischen sehr guten Evaluationen in Spezialvorlesungen und auch sehr guten Evaluationsergebnissen in großen Vorlesungen abwägen. Das Komitee hat sich entschieden, den Lehrpreis dieses Jahr für besondere Leistungen in einer großen Einführungsvorlesung zu vergeben. Von den Studierenden wurde deutlich hervorgehoben, dass die Struktur der Vorlesung sehr schlüssig und die Erklärungen hilfreich waren. So kommentiert Fachschaftler Friedrich Bach: „Das Lehrpreiskomitee war überrascht, dass die Studierenden die Vorlesung so gut bewertet haben, obwohl es für viele Studierende eine der schwersten Vorlesungen im Studium ist.“

\begin{aufgabe}[Erzeugt folgende Werte mit \texttt{siunitx}:]
	\item \SI{12345,67890}{\ohm}
	\item \num{1 +- 2i}
	\item \SI{.3e45}{\m\cubed}
	\item \num{1.654 x 2.34 x 8,99}
	\item \SI{1,793 +- 0,026}{\micro\F}
	\WarningFilter{latexfont}{Font shape}
	\item \SI{9,81+-0,52}{\m\per\s\squared}
	\DeactivateWarningFilters
\end{aufgabe}

\begin{aufgabe*}
	Zitiert den Nolting 4\footcite{nolting4} und den Demtröder 1\footcite{demtröder1} im Dokument und erstellt das Literaturverzeichnis (\lstinline!\printbibliography!)!
\end{aufgabe*}

\begin{aufgabe*}
Baut das Dokument in der Datei \texttt{tableofcontents.pdf} nach!

\emph{Hinweis:} Das Inhaltsverzeichnis wird mit dem Befehl \lstinline!\tableofcontents! erzeugt.
\end{aufgabe*}

\begin{aufgabe*}
In der Datei \texttt{beispielprotokoll.pdf} findet ihr ein Protokoll zu einem der Versuche im Grundpraktikum. Baut es so originalgetreu wie möglich nach!

Die nötigen Bilder findet ihr im Ordner \texttt{beispielprotokoll\_bilder}.
\end{aufgabe*}

\begin{aufgabe*}
	Erzeugt die Datei \texttt{prony.pdf} mit gnuplot!
	
	\enlargethispage{\baselineskip}
	\emph{Hinweise}:
	\begin{itemize}
		\item Die Messwerte befinden sich in der Datei \texttt{prony.csv}.
		\item Ihr dürft konstante Werte für die Fehlerbalken auf der $η$-Achse wählen!
		\item Nutzt die folgenden Formeln und Werte, um den Wirkungsgrad $η$ zu berechnen:
		\begin{align*}
			P_\text{el} &= \num{14,6} ⋅ \num{0,048} ⋅ \frac{1000}{3} \si{\W} \\
			r &= \SI{0,5}{\m} \\
			P &= 2πfFr \\
			η &= \frac{P}{P_\text{el}}
		\end{align*}
		\item Die Ausgleichskurve ist eine Parabel.
	\end{itemize}
\end{aufgabe*}


\end{document}
